\documentclass[12pt]{article}
\usepackage{amsmath,amssymb,amsfonts}
\usepackage[margin=1in]{geometry}

\begin{document}

\section*{Chapter 7, Problem 20}

\textbf{Problem:} Using the eigenfunction expansion method, show that the potential energy in a two-dimensional, charge-free region bounded by circles of radii $R_1$ and $R_2$ is given by
\[
V(r,\theta) = \oint G_1(r,\theta;\,R_1,\theta')\frac{\partial V}{\partial r'}\bigg|_{r'=R_1}R_1\,d\theta' + \oint G_2(r,\theta;\,R_2,\theta')\frac{\partial V}{\partial r'}\bigg|_{r'=R_2}R_2\,d\theta',
\]
where $V(R_1,\theta)$ is the known value of the potential on the first surface as a function of $\theta$, and $V(R_2,\theta')$ is its value on the second surface. Find the eigenfunction expansion of $G_1$ and $G_2$, using the fact that $\nabla^2 V = 0$ is an Hermitian operator if one imposes periodic boundary conditions on $(0,2\pi)$. Why do we require such conditions? Does the result you expect? Show that $G_1(R_1,\theta;R_1,\theta') = \delta(\theta-\theta')$ and $G_1(R_2,\theta;R_2,\theta') = 0$, and give the result you expect.

\bigskip
\textbf{Solution:}

In the annular region $R_1 \le r \le R_2$, Laplace's equation in polar coordinates is:
\[
\nabla^2 V = \frac{1}{r}\frac{\partial}{\partial r}\Bigl(r\frac{\partial V}{\partial r}\Bigr) + \frac{1}{r^2}\frac{\partial^2 V}{\partial\theta^2} = 0.
\]

\textbf{Periodic BCs:} We require $V(\theta + 2\pi) = V(\theta)$ since $\theta$ is an angular variable. This makes the angular part of $\nabla^2$ Hermitian.

\textbf{Separation of variables:} Write $V(r,\theta) = R(r)\Theta(\theta)$. The angular equation gives $\Theta'' = -n^2\Theta$ with solutions $e^{in\theta}$, $n = 0, \pm 1, \pm 2, \ldots$

The radial equation for $n \neq 0$ gives $R(r) = A_n r^n + B_n r^{-n}$, and for $n = 0$: $R = A_0 + B_0\ln r$.

\textbf{Green's function construction:} The Green's function for Laplace's equation in the annular region satisfies:
\[
\nabla^2 G = \frac{1}{r}\delta(r-r')\delta(\theta-\theta')
\]
with $G = 0$ on both boundaries $r = R_1$ and $r = R_2$.

For each angular mode $n$, the radial Green's function is constructed from two solutions: one vanishing at $R_1$ and one vanishing at $R_2$.

For $n \neq 0$, the solution vanishing at $R_1$ is $\phi_1(r) = r^n - R_1^{2n}r^{-n}$ and the solution vanishing at $R_2$ is $\phi_2(r) = r^n - R_2^{2n}r^{-n}$.

The Green's function is:
\[
G(r,\theta;r',\theta') = \frac{1}{2\pi}\sum_{n=-\infty}^{\infty}g_n(r,r')\,e^{in(\theta-\theta')},
\]
where $g_n(r,r')$ is the radial Green's function for mode $n$, constructed from $\phi_1$ and $\phi_2$ in the standard way.

The boundary Green's functions $G_1$ and $G_2$ are obtained via Green's theorem. For the Dirichlet problem where $V$ is specified on the boundaries:
\[
V(r,\theta) = -\oint_{R_1} V(R_1,\theta')\frac{\partial G}{\partial r'}\bigg|_{r'=R_1}R_1\,d\theta' - \oint_{R_2}V(R_2,\theta')\frac{\partial G}{\partial r'}\bigg|_{r'=R_2}R_2\,d\theta'.
\]

We identify:
\[
G_1(r,\theta;R_1,\theta') = -R_1\frac{\partial G}{\partial r'}\bigg|_{r'=R_1}, \quad G_2(r,\theta;R_2,\theta') = -R_2\frac{\partial G}{\partial r'}\bigg|_{r'=R_2}.
\]

\textbf{Verification:} At $r = R_1$: $G(R_1,\theta;r',\theta') = 0$ for all $r'$ in the interior. But $\partial G/\partial r'|_{r'=R_1}$ contains the delta function information. Specifically:
\[
G_1(R_1,\theta;R_1,\theta') = \delta(\theta - \theta'),
\]
which ensures that $V(R_1,\theta) = \int V(R_1,\theta')\delta(\theta-\theta')\,d\theta' = V(R_1,\theta)$. $\checkmark$

Similarly, $G_1(R_2,\theta;R_1,\theta') = 0$ because the Green's function vanishes at $r = R_2$, so specifying $V$ on the inner boundary does not directly contribute at the outer boundary (the contribution comes through $G_2$ instead). $\blacksquare$

\end{document}
